\documentclass[5p,times]{elsarticle}

\usepackage{amsmath,amssymb,amsfonts}
\usepackage{graphicx}
\usepackage{booktabs}
\usepackage{bm}
\usepackage[hidelinks]{hyperref}

\renewcommand{\thetable}{\Roman{table}}
\newcommand{\sgn}{\mathrm{sgn}}
\newcommand{\E}{\mathbf{E}}
\newcommand{\Pm}{P_m}
\newcommand{\Pe}{P_e}
% ---- numbers exported from recompute_all.py (single source of truth) ----
\newcommand{\rmseFault}{2.3}          % pooled stable-trajectory RMSE, fault-on region (deg)
\newcommand{\rmsePost}{36.9}          % pooled stable-trajectory RMSE, post-fault region (deg)
\newcommand{\rmseAll}{36.3}           % pooled stable-trajectory RMSE, full horizon (deg)
\newcommand{\pctACC}{95.4\%}          % main 39-bus classification accuracy (max-over-time)
\newcommand{\pctACCb}{92.9\%}         % 118-bus classification accuracy
\newcommand{\ccterr}{25.0}            % main 39-bus mean CCT error (ms)
\newcommand{\pairq}{148}              % certified pairwise bound q_pw (deg, alpha=0.1)
\newcommand{\accSeed}{94.0\%~$\pm$~1.2\%} % four-seed classification accuracy
\newcommand{\confmat}{318, 6, 38, 598} % TP, FP, FN, TN
\newcommand{\fone}{0.935}             % F1 of the stable class
\newcommand{\cctbusref}{218, 211, 185, 208, 453, 217, 196, 190, 197, 183, 234, 129}
\newcommand{\cctbuspred}{211, 195, 174, 201, 248, 206, 195, 203, 195, 186, 244, 111}
\newcommand{\cctbuserr}{$-7$, $-16$, $-12$, $-7$, $-205$, $-11$, $-1$, $+13$, $-2$, $+3$, $+10$, $-18$}
\newcommand{\randRMSE}{110.1}         % uniform-random predictor RMSE on stable test trajectories (deg)
\newcommand{\sevMLPacc}{93.0\%}       % severity-MLP classification accuracy (same encoding as operator)
\newcommand{\sevMLPcct}{47.0}         % severity-MLP mean CCT error (ms)
\newcommand{\mlpcct}{20.7}            % pointwise MLP mean CCT error under same protocol (ms)
\newcommand{\grucct}{33.5}            % GRU mean CCT error under same protocol (ms)
\newcommand{\paircov}{0.893}          % pairwise-score empirical coverage (alpha=0.1)
\newcommand{\qhatdeg}{110.1}          % per-machine band half-width q (deg, alpha=0.1)
\newcommand{\decfrac}{0\%}            % fraction of predicted-stable resolved by spread+q<pi
\newcommand{\wrrate}{1.8\%}           % wrong-release rate of the predicted-stable gate
\newcommand{\wrnum}{2}                % wrong-release count
\newcommand{\wrdem}{109}              % denominator of wrong-release rate
\newcommand{\gatepurity}{96.6\%}      % gate purity (evaluation half, protocol-correct)
\newcommand{\mondA}{76.1}             % Mondrian band half-width, peak <45 deg (deg)
\newcommand{\mondC}{124.6}            % Mondrian band half-width, peak >90 deg (deg)
\newcommand{\trajbus}{11}             % Fig 2 scenario fault bus
\newcommand{\trajtc}{0.170}          % Fig 2 scenario clearing time (s)
\newcommand{\trajband}{70.8}          % Fig 2 first-swing band half-width (deg)
\newcommand{\gatecovAll}{0.914}       % coverage over predicted-stable evaluation scenarios
\newcommand{\gatecovTS}{0.946}        % oracle-conditioned coverage over true-stable&pred-stable
\newcommand{\gatenum}{56}             % n of true-stable&pred-stable subset
\newcommand{\timCPU}{0.131}           % operator per-scenario CPU throughput (ms)
\newcommand{\timRK}{28.0}             % RK4 per-scenario CPU time (ms)
\newcommand{\timSU}{214}              % throughput speedup
\newcommand{\latEnd}{30.6}            % end-to-end single-fault CCT latency, operator+encoding (ms)
\newcommand{\latRK}{9735}             % end-to-end single-fault CCT latency, TDS (ms)
\newcommand{\latSU}{318}              % end-to-end speedup
\newcommand{\latSweep}{12.3}          % 100-point sweep end-to-end (ms)
\newcommand{\cctNtwo}{39.7}             % N-2 sampled-pairs mean CCT error (ms)

\begin{document}

\begin{frontmatter}

\title{A Physics-Encoded Neural Operator with Conformal Trajectory Bands for Fast Transient Stability Screening}

\author[1]{Zhenyu~Liu\corref{cor1}}
\ead{3293857014@qq.com}
\author[1]{Jiale~Wang}
\author[1]{Jiayong~Liu}
\author[1]{Zhicheng~Zeng}
\cortext[cor1]{Corresponding author.}
\affiliation[1]{organization={International Energy College, Jinan University},
            city={Guangzhou},
            postcode={510632},
            country={China}}

\begin{abstract}
Transient stability assessment (TSA) predicts whether a power system remains synchronized
after a large disturbance and estimates the critical clearing time (CCT) of faults. Time-domain
simulation (TDS) is accurate but slow, whereas machine-learning classifiers are fast but discard
the rotor-angle trajectory and provide no error bound. We introduce a physics-encoded neural
operator mapping a fault scenario to the full center-of-inertia rotor-angle trajectory, with a
distribution-free, finite-sample error band attached through split conformal prediction. The key
enabler is a physics-informed input encoding: a fault is represented by the vector of initial
accelerations it induces on each machine, computed from the swing equation rather than a bus
index, letting the operator generalize to unseen fault locations and $N{-}1$/$N{-}2$ topologies.
On the IEEE 39-bus system it classifies stability with \pctACC{} accuracy under a max-over-time
$\pi$-separation criterion, attains a mean absolute CCT error of \ccterr{} ms, and runs
$\approx\timSU{}\times$ faster than TDS on identical hardware. A directly calibrated pairwise
score bounds the machine-angle spread error to \pairq{}$^\circ$ at $\alpha{=}0.1$; the band is
coverage-valid, but the screening rule then resolves none of the scenarios, so the resolvable
fraction is reported across $\alpha$ with a TDS fallback---the conformal output is a calibrated
trajectory band, not a stability certificate. Retrained on the IEEE 118-bus system it reaches
\pctACCb{} accuracy, and the encoding carries over unchanged to higher-fidelity generator
models. A physics-residual loss does not help, and generalization across operating points remains
open. To our knowledge, this is the first work coupling a trajectory-level operator with
conformal trajectory bands and CCT extraction in TSA.
\end{abstract}

\begin{keyword}
conformal prediction \sep critical clearing time \sep neural operator \sep physics-informed
machine learning \sep power system dynamics \sep transient stability
\end{keyword}

\end{frontmatter}

\section{Introduction}
Transient stability is the ability of a power system to remain in synchronism
after a large disturbance such as a three-phase fault. Its assessment is a safety-critical,
recurring task: for a given contingency and fault clearing time, system operators must know
whether the system survives and, ideally, how much clearing-time margin remains before it loses
synchronism~\cite{sauer1998dynamics, kundur1994power}. The standard tool is time-domain
simulation (TDS), which integrates the nonlinear swing equations and is accurate but
computationally expensive, particularly when screening many fault locations, clearing times, and
topology changes~\cite{stott1979power}.

This cost has motivated a large body of machine-learning surrogates for TSA: classifiers such as
support vector machines, decision trees, or deep networks~\cite{alimi2020review} that map
operating-condition or fault features to a binary stable/unstable label. Fast as they are, these
surrogates discard the rotor-angle trajectory (and thus any stability \emph{margin} or CCT
estimate), do not generalize naturally across fault locations (a one-hot bus index yields no
signal for unseen buses), and output no uncertainty or error bound, so an operator cannot know
when the surrogate is trustworthy.

We address all three limitations with a physics-encoded neural operator equipped with conformal
trajectory bands. The operator is a Deep Operator Network (DeepONet)~\cite{lu2021learning}, whose
branch encodes a fault scenario and whose Fourier-feature trunk encodes time; their contraction
yields the full center-of-inertia (COI) rotor-angle trajectory rather than merely a label. The
fault scenario is encoded by a \emph{physics-informed} vector: the initial acceleration
$(P_{mi}-P_{ei}^{\mathrm{f}})/M_i$ that the fault induces on each machine, computed directly from
the swing equation, which allows the operator to generalize across unseen fault locations. Split
conformal prediction~\cite{lei2018distribution} attaches a finite-sample, distribution-free error
band to every predicted trajectory, together with a tighter score that directly bounds the
pairwise angle difference entering the stability criterion. Throughout, ``physics-informed''
refers to the input encoding derived from the swing equation, not to a residual-constrained model,
and the conformal output is a calibrated trajectory band, not a stability certificate.

Our contributions are as follows.
\begin{itemize}
\item We couple, to our knowledge for the first time, a trajectory-level neural operator with
conformal trajectory bands for \emph{rotor-angle} trajectories and with CCT extraction. Prior
neural operators for power systems target component-level dynamics~\cite{karampinis2025neural};
conformalized operator regression has appeared for post-fault bus-voltage
trajectories~\cite{mollaali2025conformalized}, targeting voltage rather than rotor-angle stability
and without CCT extraction; and prior conformal work on TSA predicts only a binary stability
label~\cite{lu2026confidence}. Section~\ref{sec:related} positions these works relative to
ours, and the ``first'' claim is scoped to rotor-angle trajectory bands with CCT extraction.
\item We show that a physics-informed \emph{input} encoding---the swing-equation acceleration
vector---enables generalization across unseen fault locations, and we quantify the difference
against one-hot and graph baselines.
\item We provide a split-conformal calibration layer that yields a finite-sample,
distribution-free error band on the trajectory and a directly calibrated pairwise-angle bound
below the $\pi$ threshold at $\alpha{=}0.1$, together with the resolvable fraction, the
wrong-release rate, and a CCT estimate compared against classical direct criteria; the derived
CCT band is a heuristic interval without a formal guarantee (Section~\ref{sec:method}).
\item We extend the encoding with a post-fault severity vector and show that the operator
generalizes to held-out $N{-}1$ line trips and $N{-}2$ line pairs.
\item We report trajectory accuracy together with an ablation of an auxiliary physics-residual
loss, including a negative result for this system and protocol: physics is injected more reliably
through the input encoding than through a residual loss.
\end{itemize}

\section{Related Work}
\label{sec:related}

\subsection{Neural operators and physics-informed learning in power systems}
Neural operators learn maps between function spaces rather than between fixed-size vectors; the
dominant architectures are DeepONet~\cite{lu2021learning} (branch--trunk contraction) and the
Fourier Neural Operator~\cite{li2021fourier}. Physics-informed neural networks
(PINNs)~\cite{raissi2019physics} embed governing-equation residuals into the training loss.

Recent work has begun applying these tools to power systems. Karampinis et
al.~\cite{karampinis2025neural} propose a DeepONet-based framework for modeling power-system
\emph{components}, targeting fast dynamic simulation and digital twins, but do not address TSA,
the CCT, or uncertainty quantification. Moya et al.~\cite{moya2025conformal} attach
split-conformal intervals to DeepONets (Conformalized-DeepONet) for operator regression on
canonical dynamical systems, and physics-informed graphical surrogates have been applied to power
flow and optimal power flow~\cite{huang2025optimal, chen2025power}. Our work instead targets the
post-fault rotor-angle trajectory operator and couples it with calibrated trajectory
uncertainty quantification.

\subsection{Learning-based transient stability assessment}
Learning-based TSA has long been dominated by classifiers that predict a stable/unstable label
from hand-crafted features~\cite{alimi2020review}; more recent graph models with
distribution-aware pooling~\cite{chen2022distribution} report high accuracies on the IEEE 39-bus
system under in-sample or per-sample protocols. Under our held-out-bus protocol with the same raw
features, however, a standard graph convolutional classifier reaches only $76.0\%$ and a stronger
residual graph model $77.3\%$ (Section~\ref{sec:results}), and graph models produce no
trajectory, CCT, or uncertainty bound.

A few works predict the post-fault trajectory rather than only the label; these are closest in
spirit to ours, but they provide no error bounds. Event-structured PINNs model the pre-fault,
fault-on, and post-fault stages with hard state chaining and a differentiable CCT
boundary~\cite{hao2026event}; their error-propagation theorem is complementary to our
finite-sample conformal band. PINN-based acceleration of time-domain simulation
\cite{stiasny2021transient} trains a network with physics residuals, whereas we learn the
scenario-to-trajectory map with a physics-informed input encoding; fast CCT calculation for
mixed synchronous/asynchronous generation~\cite{wang2025fast} is complementary to our TDS-based
reference.

Confidence-aware TSA has begun to emerge. Lu et al.~\cite{lu2026confidence} combine a large
language model with conformal prediction to produce prediction sets for a binary stability label,
but they predict neither the trajectory nor the CCT and do not use a physics-informed operator.
The closest works to ours are the conformalized operator regression of Mollaali et
al.~\cite{mollaali2025conformalized}, which predicts post-fault \emph{bus-voltage} trajectories
without rotor angles, CCT extraction, or topology generalization, and the Conformalized-DeepONet
framework of Moya et al.~\cite{moya2025conformal}, which establishes operator-regression bands on
canonical dynamical systems. Probabilistic and Bayesian operator frameworks
(DeepONet-grid-UQ)~\cite{moya2023deeponet} assume Gaussian or posterior-sampling uncertainty,
and risk-controlling quantile operators (UQNO)~\cite{ma2024calibrated} learn a second operator
for pointwise radii; our split-conformal bands are distribution-free and require only a
calibration set on top of the trained operator.

Across the axes that matter for TSA---physics-derived input encoding, unseen fault-bus and
topology generalization, the calibration unit, and CCT extraction---no single prior work combines
a physics encoding that generalizes across fault locations and topologies, a directly calibrated
bound on the pairwise spread entering the stability criterion, and CCT extraction.

Classical direct methods remain strong CCT baselines: the transient-energy-function (TEF) / potential-energy-boundary-surface (PEBS)
method~\cite{athay1979transient} and the individual-machine equal-area criterion
(IMEAC)~\cite{wang2018individual} extract the CCT from machine-level energy criteria
(Section~\ref{sec:results}). Our contribution is the combination of a trajectory-level neural
operator with a physics-informed fault encoding and a conformal error band, which yields a CCT
estimate alongside the trajectory; the derived CCT band is heuristic and carries no formal
coverage guarantee (Section~\ref{sec:method}).

\section{Problem Formulation}
\label{sec:problem}

\subsection{Classical multi-machine model}
We consider the classical model of transient stability, in which each synchronous generator is
represented as a constant voltage source $E_i$ behind a transient reactance
$x'_{di}$~\cite{sauer1998dynamics, kundur1994power}. Under the standard reduction, loads are
converted to constant admittances and the network is Kron-reduced to the internal generator
nodes, yielding an $n\times n$ admittance matrix $\bm{Y}$. The dynamics of generator $i$ obey the
swing equations
\begin{align}
\frac{d\delta_i}{dt} &= \omega_s(\omega_i - 1), \label{eq:swing1}\\
\frac{d\omega_i}{dt} &= \frac{P_{mi} - P_{ei}(\bm{\delta}) - D_i(\omega_i - 1)}{2H_i}, \label{eq:swing2}
\end{align}
where $\delta_i$ is the rotor angle, $\omega_i$ the per-unit speed, $H_i$ the inertia constant,
$D_i$ the damping coefficient, $P_{mi}$ the mechanical power, and $\omega_s$ the synchronous
speed. The electrical power is
\begin{equation}
P_{ei}(\bm{\delta}) = \sum_{j=1}^{n} E_i E_j\bigl[ G_{ij}\cos(\delta_i-\delta_j)
+ B_{ij}\sin(\delta_i-\delta_j)\bigr],
\end{equation}
with $G_{ij}=\Re\{Y_{ij}\}$ and $B_{ij}=\Im\{Y_{ij}\}$. We use the undamped case $D_i=0$, which
is standard for the classical model and gives the undamped (energy-margin) stability boundary; a
damped variant is examined in Section~\ref{sec:results}. Because the
equations depend only on angle differences, it is convenient to work in the center-of-inertia
(COI) frame, defining
\begin{equation}
\tilde\delta_i = \delta_i - \delta_{\mathrm{COI}}, \qquad
\delta_{\mathrm{COI}} = \frac{\sum_i M_i \delta_i}{\sum_i M_i}, \quad M_i = 2H_i.
\end{equation}

A three-phase fault at a bus is modeled by adding a large shunt admittance to that bus and
reducing the network during the fault; the fault is cleared after a clearing time $t_c$,
optionally tripping a line (an $N{-}1$ contingency). Reference trajectories integrate
\eqref{eq:swing1}--\eqref{eq:swing2} with a fourth-order Runge--Kutta scheme at a $0.5$ ms step,
splitting each step at the clearing instant $t_c$ so the fault-on network is used exactly until
$t_c$ (event alignment); the state is recorded on a $0.01$ s output grid. A trajectory is
\emph{stable} if the maximum pairwise angle difference remains below $\pi$ at every output-grid
point over the simulation horizon---the $\pi$-separation criterion, the operational label used
throughout the paper. It applies to the selected model, horizon, and criterion and is not claimed
to coincide with infinite-horizon synchronism for arbitrary systems
(Section~\ref{sec:discussion}).

\subsection{Problem statement}
Let a fault scenario be described by a fault location (bus), an optional set of tripped lines, and
a clearing time $t_c$. We seek the map
$\mathcal{G}\colon(\text{scenario})\mapsto\{\tilde\delta_i(t)\}_{i=1}^{n}$, $t\in[0,T]$,
returning the COI-referenced rotor-angle trajectory, together with (i) a stability
classification under the $\pi$-separation criterion, (ii) the critical clearing time
$t_{\mathrm{CCT}}$, the first clearing time at which the criterion is violated as $t_c$
increases from zero---a dense sweep found a single dominant stable-to-unstable crossing on most
buses (1--5 local flips at the $5$ ms resolution, concentrated on the two largest-CCT buses), so
on the held-out buses this coincides with the largest stable $t_c$---and (iii) a calibrated error
band on the trajectory, in the statistical sense made precise in Section~\ref{sec:method}.

\section{Method}
\label{sec:method}

\subsection{Physics-informed fault encoding}
A central obstacle to generalization across fault locations is that a bus index is an arbitrary
label: a one-hot vector gives the branch network no information about physically similar buses. We instead encode a fault at bus $k$ by the vector of
initial accelerations it induces, computed from the swing equation at the pre-fault equilibrium,
\begin{equation}
a_i(k) = \frac{P_{mi} - P_{ei}^{\mathrm{f}}(k)}{M_i},
\label{eq:severity}
\end{equation}
where $P_{ei}^{\mathrm{f}}(k)$ is the electrical power of machine $i$ during the fault at bus
$k$, evaluated at the pre-fault angles---precisely the initial rate of change of per-unit speed in
\eqref{eq:swing2} (the angular acceleration is $\omega_s a_i$). Dividing by $M_i$ makes machines
of widely different inertia contribute comparably, and continuity in the fault admittance maps
physically similar faults to nearby encodings. The branch input is the concatenation of
$\bm{a}(k)\in\mathbb{R}^n$ and the normalized clearing time. The encoding is not claimed to be
injective---two scenarios with similar encodings can still differ in later response near the
stability boundary---but Section~\ref{sec:results} shows it transfers the information needed for
this task, while the conformal band flags the scenarios where the prediction is least reliable.

For $N{-}1$ contingencies, in which a line is tripped at fault clearing, the post-fault reduced
admittance and the post-fault dynamics both change. We encode this change symmetrically with a
\emph{post-fault} severity vector,
\begin{equation}
\tilde a_i(l) = \frac{P_{mi} - P_{ei}^{\mathrm{post}}(l)}{M_i},
\label{eq:postseverity}
\end{equation}
where $P_{ei}^{\mathrm{post}}(l)$ is the post-fault electrical power at the pre-fault angles with
line $l$ removed; for the intact network $\tilde a_i=0$. The full scenario encoding concatenates
the fault-on severity \eqref{eq:severity}, the post-fault severity \eqref{eq:postseverity}, and
the clearing time. For $N{-}2$ contingencies the same vector is evaluated with both lines removed:
$\tilde{\bm a}$ is always an $n$-dimensional vector evaluated at the pre-fault angles, so the
encoding dimension is unchanged and the same branch network applies without modification.

\subsection{Neural operator architecture}
We use a DeepONet. The branch network $b_\theta$ maps the scenario encoding to a matrix of basis
coefficients $\bm{B}\in\mathbb{R}^{n\times p}$; the trunk network $\tau_\phi$ maps time to $p$
basis functions. The predicted trajectory is the contraction
\begin{equation}
\hat{\bm{\delta}}(s,t) = \bm{B}(s)\,\tau(t) \in \mathbb{R}^n .
\end{equation}
Because the swing dynamics are oscillatory, we feed the trunk a Fourier feature embedding of
time, $\{\cos(2\pi k t/T),\sin(2\pi k t/T)\}_{k=1}^{K}$, alongside the normalized raw time,
before a small multilayer perceptron,
so the trunk represents the natural oscillation modes more efficiently than a plain perceptron.
The trunk is smooth
in $t$, so its time derivatives are available by automatic differentiation for the optional
physics residual below. Fig.~\ref{fig:arch} summarizes the proposed pipeline.

\begin{figure}[t]
\centering
\includegraphics[width=\columnwidth]{fig_arch.pdf}
\caption{Overview of the proposed physics-informed neural operator ($N{-}0$ configuration shown;
$N{-}1$/$N{-}2$ additionally use the post-fault severity vector): a fault scenario is mapped
through the physics-informed encoding to a branch network, whose coefficients contract with a
Fourier-feature trunk to yield the predicted rotor-angle trajectory $\hat{\bm\delta}$; split
conformal prediction then attaches a distribution-free error envelope.}
\label{fig:arch}
\end{figure}

\subsection{Physics-residual loss (auxiliary)}
Following the PINN paradigm, we optionally add a residual of the post-fault swing equation in the
COI frame. Differentiating the COI definition under $D_i=0$ gives
$M_i \ddot{\tilde\delta}_i = \omega_s(P_{mi}-P_{ei}) - M_i\omega_s c(t)$, where
$c(t)=\sum_j(P_{mj}-P_{ej})/M_{\mathrm{tot}}$ is the COI acceleration and
$M_{\mathrm{tot}}=\sum_j M_j$. Dividing by $M_i$ so that machines of widely different inertia
contribute comparably (otherwise the $H{=}500$ s aggregate machine at bus~39 dominates the sum),
the per-machine residual enforced at a set $\mathcal{C}$ of collocation times in the post-fault
region $t\ge t_c$ is
\begin{equation}
\mathcal{L}_{\mathrm{phys}} = \frac{1}{|\mathcal{C}|}\sum_{t\in\mathcal{C},\,i}
\Bigl[ \ddot{\tilde\delta}_i - \omega_s(P_{mi}-P_{ei}(\tilde{\bm{\delta}}))/M_i + \omega_s c(t) \Bigr]^2 .
\label{eq:physloss}
\end{equation}
The per-machine division by $M_i$ makes the term comparable across machines; it retains
acceleration units and is a soft penalty rather than a hard constraint. The $\omega_s$ factor
enters through the frequency-to-angle relation and fixes the scaling between the curvature and
power-mismatch terms. The total loss is the data loss plus $\lambda\mathcal{L}_{\mathrm{phys}}$,
enforced only in the post-fault region (the fault-on admittance varies per scenario) at every
fifth output time point; denser collocation is ablated in Section~\ref{sec:results}.

\subsection{Conformal calibration}
The trajectory operator alone gives no indication of its own error. We use split conformal
prediction~\cite{lei2018distribution} to attach a finite-sample, distribution-free band to the
single-machine COI angles, calibrated on \emph{stable} scenarios: unstable trajectories diverge
and carry errors orders of magnitude beyond any useful bound, so the band is defined on the stable
subpopulation by design. For a scenario $s$, the nonconformity score is the worst-case prediction error over
machines and output-grid time points,
\begin{equation}
R(s) = \max_{i,t\in\mathcal{T}_{\mathrm{grid}}} \bigl| \hat{\tilde\delta}_i(t;s) - \tilde\delta_i(t;s) \bigr| .
\end{equation}
We split the $N$ stable scenarios of the held-out test set (Section~\ref{sec:setup})---which the
operator never saw during training and which are drawn at random from the same held-out-bus
sampling distribution, so calibration and evaluation scores are exchangeable---into a calibration
half ($N_{\mathrm{cal}}$) and a disjoint evaluation half ($N_{\mathrm{eva}}$). For target
miscoverage $\alpha$, the bound $\hat q$ is the $k$-th order statistic of the calibration scores
with $k=\lceil(N_{\mathrm{cal}}+1)(1-\alpha)\rceil$, the standard finite-sample split-conformal
rule. For a new scenario exchangeable with the calibration scenarios and drawn from the same
distribution \emph{conditionally on being truly stable}, the whole predicted trajectory is then
covered simultaneously on the output grid: $\Pr\{ \max_{i,t\in\mathcal{T}_{\mathrm{grid}}}
\bigl| \hat{\tilde\delta}_i(t) - \tilde\delta_i(t) \bigr| \le \hat q \mid \text{scenario stable} \}
\ge 1-\alpha$, i.e., with probability at least $1-\alpha$ over exchangeable stable scenarios,
every single-machine COI angle at every \emph{output-grid} point lies within $\hat q$ of the
truth. Three limits are
explicit: the guarantee is marginal over the stable subpopulation (it does not cover misclassified
unstable scenarios); it applies to the finite output grid ([Rev. R2] a $2$ ms dense-grid check on $200$ held-out
scenarios found no missed criterion crossings, i.e., no label flips relative to the $10$ ms
output grid); and it is a calibrated trajectory band, not a
stability certificate. By the triangle inequality, the pairwise angle difference relevant to the
stability criterion is bounded by $2\hat q$, which at $\alpha{=}0.1$ exceeds the $\pi$ threshold.

To bound the decision quantity directly, we also calibrate a pairwise-difference score over the
signed per-machine errors,
\begin{equation}
R_{\mathrm{pw}}(s) = \max_{i,j,t\in\mathcal{T}_{\mathrm{grid}}} \bigl| \hat e_i(t) - \hat e_j(t) \bigr|, \qquad
\hat e_i(t) = \hat{\tilde\delta}_i(t) - \tilde\delta_i(t),
\end{equation}
which scores exactly the pairwise angle spread appearing in the $\pi$-separation criterion; the
errors $\hat e_i$ keep their sign, so that oppositely signed machine errors add rather than
cancel. As shown in Section~\ref{sec:results}, $R_{\mathrm{pw}}$ bounds the pairwise-difference
prediction error to \pairq{}$^\circ$ at $\alpha{=}0.1$---below the $\pi$ threshold, but too wide to certify decisions
in practice: the rule $\mathrm{spread}+\hat q_{\mathrm{pw}}<\pi$ resolves \decfrac{} of the
predicted-stable cases at $\alpha{=}0.1$, rising to $31\%$ at $\alpha{=}0.3$ and $67\%$ at
$\alpha{=}0.5$ as the bound tightens. The pairwise bound is therefore a screening aid used with a
TDS fallback for the unresolved fraction, not a stand-alone decision rule. Multiplicity-corrected
per-machine bounds (Bonferroni/\v{S}id\'ak) do not tighten the pairwise bound
(Section~\ref{sec:results}).

A CCT band can be derived from the per-machine interval: clearing times whose predicted
trajectory, shifted by $\pm 2\hat q$ in spread space, cannot reach the $\pi$ threshold are
declared stable with margin, and clearing times whose shifted spread exceeds $\pi$ for \emph{all}
shifts are declared unstable; the remainder forms the band, found by the same binary search with
the shifted criterion. Because the trajectory bound at one $t_c$ does not transfer across $t_c$
(the CCT is a nonlinear functional of the trajectory), this CCT band is \emph{heuristic}: it
carries no formal coverage guarantee and is reported as a screening aid
(Section~\ref{sec:results}).

\subsection{Critical clearing time estimation}
Given the operator, the CCT of a fault is the first clearing time at which the criterion is
violated. Because neither the operator's criterion nor the reference is guaranteed monotone in
$t_c$, a $5$ ms sweep over $[0.01,0.8]$ s locates the first stable-to-unstable sign change, which
is then refined by local bisection to $0.1$ ms; the same first-crossing protocol (with the
system-dependent range $[0.005,0.2]$ s) is applied to the TDS reference, so the two CCT estimates
are directly comparable. The refinement tolerance is the search resolution, not the physical CCT
accuracy; the reference uses the event-aligned RK4 scheme of Section~\ref{sec:problem}
(convergence verified at a $0.25$ ms step), and both judge stability on the $0.01$ s grid.

\section{Experimental Setup}
\label{sec:setup}

\subsection{Test system and data}
We use the IEEE 39-bus New England system with $10$ machines, whose static data are the standard
published values and whose dynamic parameters ($H$, $x'_{d}$) are the canonical
classical-model values~\cite{sauer1998dynamics, athay1979transient}; the specific parameter
values and their source are tabulated in the code repository. The reference simulator integrates
\eqref{eq:swing1}--\eqref{eq:swing2} with the event-aligned RK4 scheme of Section~\ref{sec:problem}
at a $0.5$ ms step over a $4$ s horizon. We draw a fault bus uniformly over all $39$ buses and a
clearing time uniformly in $[0.03,0.55]$ s, producing approximately $40\%$ stable scenarios. The
dataset has $2160$ training, $540$ validation, and $960$ test scenarios; the test set is
constructed from $12$ fault buses held out entirely from training, so test performance measures
generalization to unseen fault locations, while the validation set is drawn from the training
buses. Event-aligned labels differ from the pre-alignment reference (which kept the fault on until
the first grid point after $t_c$) on exactly one of $960$ test and one of $2160$ training
scenarios, both at the stability boundary. The CCT search range $[0.01,0.8]$ s covers the observed 39-bus CCTs ($0.13$--$0.58$ s;
$0.13$--$0.45$ s on the held-out buses).

\subsection{Baselines and ablations}
We compare against the same DeepONet with (i) a one-hot fault encoding; (ii) a physics-residual
loss at several weights; (iii) a self-attention branch; and (iv) a Fourier-neural-operator trunk,
plus a pointwise multilayer perceptron (MLP) baseline mapping (scenario, time) to the angle
directly, graph convolutional networks (GCNs) that predict the stability label or the CCT from
raw node features (voltages, angles, loads, and a fault flag), and a gated-recurrent-unit (GRU)
sequence-to-sequence baseline. The CCT is additionally compared against the transient-energy
function (PEBS) method~\cite{athay1979transient} with the conductance path-integral terms of Athay
et al., and an IMEAC-style individual-machine criterion~\cite{wang2018individual}. The DeepONet variants, MLP, and GRU share the same dataset,
optimizer, class-balanced mean-squared loss (saturating beyond $\pm 3\pi$), and $600$-epoch
budget; the GCNs use binary cross-entropy (or mean squared error (MSE)) and $300$ epochs on raw observables, so the
GCN comparison measures the value of the severity encoding. A severity-MLP baseline receives
exactly the same encoding as the operator, and the pointwise MLP and GRU are equipped with the
same conformal calibration and CCT search, so all trajectory-level models are compared under one
protocol; parameter counts and training budgets are in the repository.

\subsection{Evaluation}
We report the root-mean-square error (RMSE) of the predicted trajectory over stable test
scenarios, decomposed into the fault-on ($t<t_c$) and post-fault ($t\ge t_c$) regions; the
stability classification accuracy under the max-over-time $\pi$-separation criterion, with the
confusion matrix and precision, recall, and $F_1$ of the stable class; the conformal band
half-width and its empirical coverage with counts and binomial intervals; the CCT error against
the event-aligned TDS first-crossing search; and the wall-clock speedup over the RK4 reference. ``RMSE''
refers to the \emph{pooled} RMSE over all stable scenarios, machines, and time steps in the
relevant region, computed in one pass; it is dominated by near-boundary scenarios and is not
comparable to a per-scenario time-averaged mean absolute error. [Rev. R9] The results are not artifacts of the hyperparameters: halving the Fourier harmonics
to $K{=}8$ degrades the pooled RMSE to $45.9^\circ$ ($94.3\%$ accuracy), and reducing the loss
saturation to $\pm2\pi$ yields $94.7\%$ accuracy with a slightly better pooled RMSE of
$34.3^\circ$; we retain $\pm3\pi$ because it is selected on validation data and gives more
stable fits to diverging unstable trajectories. The basis dimension $p$ is ablated in
Table~\ref{tab:ablation} and Section~\ref{sec:sampleeff}. All main-table numbers are
exported by a single evaluation script from one model and one dataset; four-seed statistics are
reported separately where available.

The operator is implemented in PyTorch ($2.13$, CUDA $12.6$) and trained on a single NVIDIA RTX
4070 GPU; CPU timings use a single thread of an AMD Ryzen 7 7700. Both networks use tanh
activations: the branch has three hidden layers of width 512 and outputs $n\times p$ basis
coefficients with $p=512$; the trunk has three hidden layers of width 128 on top of $K=16$
Fourier harmonics. We train with Adam (learning rate $10^{-3}$, cosine schedule, $600$ epochs,
batch size $128$). The main results (Table~\ref{tab:main}) are for the seed-0 model under the
max-over-time criterion; the same script also exports the confusion matrix, and across four
training seeds the classification accuracy is \accSeed{} and the per-scenario mean RMSE is
$20.2^\circ~\pm~1.4^\circ$, so the reported values are representative. All test metrics other than
conformal coverage (RMSE, classification, CCT) are computed on the full $960$-scenario held-out
set; only conformal coverage is restricted to a disjoint evaluation half
(Section~\ref{sec:method}). The speedup is reported under two separate protocols on identical CPU hardware,
single-threaded, with median/P95 statistics.
\emph{Throughput}: over repeated batches of $64$ scenarios the operator evaluates a $401$-point
trajectory at \timCPU{} ms per scenario against \timRK{} ms for the RK4 reference ($0.5$ ms step,
$4$ s horizon), a $\approx\timSU{}\times$ gap. \emph{End-to-end latency}: a single fault
scenario, including the severity encoding (Kron reduction, $\approx0.13$ ms per bus on the 39-bus
system) and the sequential $16$-iteration CCT binary search, takes \latEnd{} ms versus
\latRK{} ms for the same brute-force search with TDS, a \latSU{}$\times$ gap; a $100$-point
clearing-time sweep costs \latSweep{} ms end-to-end versus $2.3$ s for RK4. On the 118-bus system
the encoding ($\approx4.1$ ms per fault) dominates a single scenario. [Rev. R11] The encoding is
a Kron reduction whose measured cost is $0.13$ ms per fault on the 39-bus system and $4.1$ ms on
the 118-bus system---still a small fraction of one TDS run, but it grows with network size---and
the operator forward pass is linear in the number of output grid points and in the number of
machines (the branch input dimension is $n+1$). The RK4 cost scales
inversely with the reference step size, whereas the operator's cost depends only on the number of
output time points ($401$ per trajectory), so the speedup grows for finer reference steps; all
timing details are in the repository (\url{https://github.com/La0Fe1/operator-transient-stability}).

\section{Results}
\label{sec:results}

\subsection{Trajectory prediction and classification}
Table~\ref{tab:main} summarizes the main results. The physics-informed operator predicts the
fault-on segment of the trajectory to \rmseFault{}$^\circ$ pooled RMSE and the post-fault segment
to \rmsePost{}$^\circ$ (full-horizon \rmseAll{}$^\circ$); the post-fault error is concentrated
near the stability boundary---the sensitivity the conformal band is designed to capture---with
$10.5^\circ$ in the first-swing window and $41.8^\circ$ in the subsequent oscillation. The
operator classifies stability under the max-over-time $\pi$-separation criterion with \pctACC{}
accuracy and confusion matrix (TP, FP, FN, TN) $=$ \confmat{} ($F_1$ for the stable class
\fone{}), and attains a mean absolute CCT error of \ccterr{} ms (max $204.9$ ms on bus 12), at
a $\approx\timSU{}\times$ throughput speedup over the RK4 reference.

For reference, the transient-energy-function (PEBS) approach of Athay et
al.~\cite{athay1979transient} with the conductance path-integral terms attains a mean CCT error
of $18.7$ ms with a conservative bias, and an IMEAC-style individual-machine
criterion~\cite{wang2018individual} yields $9.6$ ms. The operator's \ccterr{} ms mean absolute
error is larger than both and carries a conservative bias
($-21$ ms), dominated by bus 12 (reference CCT $453$ ms, estimated $248$ ms), where the criterion
has multiple local stable-to-unstable crossings; the first-crossing search of
Section~\ref{sec:method} reproduces the previous binary-search CCTs within $3$ ms per bus.
Classical direct methods remain the more accurate CCT estimators; the operator's value is that it
additionally returns the full trajectory and a conformal band without any fault-on or post-fault
integration. Fig.~\ref{fig:trajectory} shows predicted versus true COI trajectories for a
representative near-boundary scenario.

\begin{table}[t]
\centering
\caption{Main results on the IEEE 39-bus system.}
\label{tab:main}
\begin{tabular}{lc}
\toprule
Metric & Value \\
\midrule
Stable trajectory RMSE (fault-on region) & \rmseFault{}$^\circ$ \\
Stable trajectory RMSE (post-fault region) & \rmsePost{}$^\circ$ \\
Stability classification accuracy (max-over-time) & \pctACC{} \\
\;\; precision / recall / $F_1$ & $0.981$ / $0.893$ / \fone{} \\
CCT error (mean $|$err$|$ / max) & \ccterr{} / $204.9$ ms \\
Conformal band half-width ($\hat q$, $\alpha{=}0.1$) & \qhatdeg{}$^\circ$ \\
Conformal empirical coverage & $0.882$ ($157/178$) \\
Pairwise bound ($R_{\mathrm{pw}}$, $\alpha{=}0.1$) & \pairq{}$^\circ$ \\
Speedup vs.\ RK4 (CPU, throughput) & $\approx\timSU{}\times$ \\
\bottomrule
\end{tabular}
\end{table}

\begin{figure}[t]
\centering
\includegraphics[width=\columnwidth]{fig_trajectory.pdf}
\caption{Predicted (dashed) versus true (solid) COI rotor-angle trajectories for the two most
affected machines in a representative near-boundary stable scenario (fault bus \trajbus{},
$t_c=\trajtc{}$ s); the shaded band is the first-swing conformal band ($\trajband{}^\circ$ at
$\alpha{=}0.1$). The fault-on segment ($t<t_c$, marked) is reproduced accurately, while the
post-fault oscillation carries the bulk of the error. The first-swing window is defined by the
\emph{true} trajectory (an offline diagnostic); the per-machine band does not certify the
pairwise criterion.}
\label{fig:trajectory}
\end{figure}

\subsection{Effect of the physics-informed encoding}
Table~\ref{tab:ablation} reports the ablation on the fault encoding. Replacing the
physics-informed acceleration vector with a one-hot bus index degrades the stable-trajectory
RMSE from \rmseAll{}$^\circ$ to $141.4^\circ$ and the classification accuracy from \pctACC{} to
$88.0\%$, because a one-hot encoding provides no signal for unseen fault buses; the one-hot
baseline also overfits at $p{=}512$---its RMSE exceeds the \randRMSE{}$^\circ$ attained by a
uniformly random predictor in $[-180^\circ,180^\circ]$ on the same stable test trajectories, an
actual baseline computed on the test set---and reducing its capacity does not recover meaningful
generalization. The physics-informed encoding is therefore the mechanism enabling generalization
across fault locations.

A GCN classifier operating on raw node features (voltage magnitudes and angles, loads, and a
fault flag) achieves only $76.0\%$ accuracy versus \pctACC{} for the severity-encoded operator;
a stronger graph baseline with three residual layers attains
$77.3\%\pm3.9\%$ (five seeds), so the gap is not an artifact of a weak architecture, and a
GCN regressing the CCT from the same raw features attains a mean CCT error of $48.0$ ms (max
$224.4$ ms) versus \ccterr{} ms (max $204.9$ ms). These comparisons are not
architecture-controlled, but
they show the severity encoding carries more information than raw observables. Note that the
pointwise MLP's higher classification accuracy does not carry over to trajectory quality
($43.1^\circ$ versus \rmseAll{}$^\circ$ pooled RMSE, without continuous-time evaluation;
Table~\ref{tab:ablation}). A \emph{severity-MLP}
receiving exactly the same encoding attains \sevMLPacc{} accuracy and \sevMLPcct{} ms mean CCT
error (Table~\ref{tab:ablation}): the encoding, not the DeepONet architecture, is the dominant
factor for generalization to unseen fault locations. A recurrent (GRU) sequence-to-sequence
baseline mapping the same encoding through a sequential decoder attains $57.3^\circ$ RMSE and
$93.9\%$ accuracy, worse than the operator's \rmseAll{}$^\circ$: the low-rank structure
outperforms a generic recurrent decoder.

\begin{table}[t]
\centering
\caption{Comparison of variants and baselines on the IEEE 39-bus system ($12$ held-out fault
buses, seed 0). RMSE is the pooled stable-trajectory RMSE; ``$-$'' marks a metric the variant does
not produce. The pointwise MLP and GRU are evaluated under the same criterion, conformal
calibration, and CCT search as the operator.}
\label{tab:ablation}
\begin{tabular}{lccc}
\toprule
Variant & RMSE & Acc. & CCT err. \\
\midrule
DeepONet, physics encoding (proposed) & \rmseAll{}$^\circ$ & \pctACC{} & \ccterr{} ms \\
DeepONet, one-hot encoding & $141.4^\circ$ & $88.0\%$ & $-$ \\
DeepONet $+$ physics-residual loss ($\lambda{=}0.001$) & $39.2^\circ$ & $95.6\%$ & $-$ \\
DeepONet $+$ physics-residual loss ($\lambda{=}0.01$) & $76.5^\circ$ & $89.9\%$ & $-$ \\
Pointwise MLP baseline & $43.1^\circ$ & $96.9\%$ & \mlpcct{} ms \\
GRU sequence baseline & $57.3^\circ$ & $93.9\%$ & \grucct{} ms \\
Severity-MLP classifier (same encoding) & $-$ & \sevMLPacc{} & $-$ \\
Severity-MLP CCT regressor (same encoding) & $-$ & $-$ & \sevMLPcct{} ms \\
GCN classifier (raw features) & $-$ & $76.0\%$ & $-$ \\
GCN CCT regressor (raw features) & $-$ & $-$ & $48.0$ ms \\
\bottomrule
\end{tabular}
\end{table}

\subsection{Generalization to $N{-}1$ and $N{-}2$ topologies}
We train the multi-topology operator ($p{=}512$) with the extended encoding of
Section~\ref{sec:method}. The $N{-}1$ experiment trains on $9990$ $N{-}0$/$N{-}1$ scenarios
(line trips drawn from $25$ training lines) and evaluates on two held-out axes: unseen fault
buses and $10$ unseen trip lines. The $N{-}2$ experiment trains on $8748$ scenarios that
additionally include $N{-}2$ line-pair scenarios drawn from the training lines, and evaluates on
held-out line pairs. On held-out line trips the operator
achieves $33.6^\circ$ stable-trajectory RMSE and $95.9\%$ accuracy, demonstrating that the
post-fault severity encoding lets the operator generalize to topologies never seen in training;
on held-out fault buses the RMSE is $64.3^\circ$ and the accuracy $93.5\%$ (Table~\ref{tab:n1}),
reflecting the harder combined task. The same encoding and architecture apply to $N{-}2$ contingencies (two lines tripped)
without modification: on held-out line pairs the operator attains $43.4^\circ$ RMSE and $92.9\%$
accuracy, because the post-fault severity vector is simply evaluated from the admittance matrix
with both lines removed. The multi-topology operator also estimates the CCT on held-out
combinations: $19.7$ ms mean absolute error on unseen trip lines (sampled), $34.1$ ms on held-out fault
buses (sampled), and \cctNtwo{} ms on held-out $N{-}2$ line pairs (sampled); split-conformal coverage on
these splits is reported in Section~\ref{sec:results}.

\begin{table}[t]
\centering
\caption{Generalization of the multi-topology operator (training mixes include $N{-}0$/$N{-}1$/$N{-}2$): pooled
stable-trajectory RMSE, classification accuracy, and CCT error (mean over held-out
bus--line combinations).}
\label{tab:n1}
\begin{tabular}{lccc}
\toprule
Held-out split & RMSE & Acc. & CCT err. \\
\midrule
Trip lines ($N{-}1$, unseen) & $33.6^\circ$ & $95.9\%$ & $19.4$ ms \\
Fault buses (mixed $N$-0/$N$-1) & $64.3^\circ$ & $93.5\%$ & $31.9$ ms \\
Line pairs ($N{-}2$, unseen) & $43.4^\circ$ & $92.9\%$ & \cctNtwo{} ms \\
\bottomrule
\end{tabular}
\end{table}

\subsection{Sample efficiency of the operator}
\label{sec:sampleeff}
Training the DeepONet
and the pointwise MLP baseline on $10\%$, $25\%$, $50\%$, and $100\%$ of the training scenarios
(three seeds each), the operator attains $56.1^\circ$ RMSE versus $65.2^\circ$ for the MLP at
$10\%$ of the data (a $\approx9^\circ$ advantage), $\approx8^\circ$ at $25\%$, and
$\approx2$--$3^\circ$ at $50\%$/$100\%$. The low-rank structure therefore confers a modest
sample-efficiency advantage in the low-data regime; with three seeds per fraction and a
$14.0^\circ$ across-seed spread at $10\%$ data, the advantage is indicative rather than
statistically established. 

\subsection{Extensions to higher-fidelity generator models and damping}
To verify that the approach carries over to more realistic generator representations, we retrain
the operator---architecture, encoding, and protocol unchanged, with the same $12$ held-out fault
buses---on data from higher-fidelity simulators, so only the ground-truth trajectories change.
On the one-axis (flux-decay) model (a third-order per-machine model~\cite{sauer1998dynamics},
Ch.~8) the retrained operator attains $38.6^\circ$ pooled stable-trajectory RMSE ($1.8^\circ$
fault-on, $39.2^\circ$ post-fault), $94.0\%$ classification accuracy, a mean CCT error of
$23.3$ ms, and a conformal band of $121.3^\circ$ half-width with $0.903$ coverage
($\alpha{=}0.1$). Adding a first-order automatic voltage regulator (a simplified IEEE DC1A
exciter with $K_A{=}20$, $T_A{=}0.2$ s and rate limits) yields a fourth-order model on which the
operator attains $25.4^\circ$ pooled RMSE ($2.2^\circ$ fault-on, $25.5^\circ$ post-fault),
$94.0\%$ accuracy, and conformal coverage $0.896$ with half-width $127.7^\circ$---the lowest
absolute RMSE of the three generator models, largely because voltage regulation smooths the
post-fault oscillation.

Adding uniform damping $D{=}1$ leaves the encoding unchanged (the damping term vanishes at
the pre-fault equilibrium), so the same architecture is retrained on damped data. The fault-on region improves to $1.8^\circ$ pooled RMSE, but the post-fault RMSE
rises to $46.4^\circ$ (versus \rmsePost{}$^\circ$ undamped) and classification drops to $93.4\%$
(CCT error $30.8$ ms): near-boundary damped trajectories decay slowly, so small errors in the
predicted decay rate accumulate over the $4$ s window (maximum worst-case error $638^\circ$). The
split-conformal band remains approximately valid but more seed-sensitive than in the undamped
case (mean coverage $0.89\pm0.05$ over five random calibration splits, range $0.80$--$0.95$;
the seed-0 split attains $0.80$); damping does not simplify the near-boundary prediction task.

\subsection{Conformal bands}
The split-conformal band achieves empirical coverage close to nominal at every level on the
$N_{\mathrm{eva}}{=}178$ evaluation scenarios: $91.6\%$ ($163/178$) at the $95\%$ target,
$88.2\%$ ($157/178$, CI $83$--$93\%$) at $90\%$, $67.4\%$ ($120/178$) at $70\%$, and $43.8\%$
($78/178$) at $50\%$ (Fig.~\ref{fig:conformal}); the same calibration on the multi-topology
operators yields coverage of $0.915$ on both held-out $N{-}1$ splits and $0.872$/$0.852$ on
$N{-}2$ buses and line pairs (half-widths $121$--$181^\circ$). A Mondrian (grouped)
calibration~\cite{vovk2005algorithmic} by the post-fault peak of the pairwise
spread---$\max_{t\ge t_c}\bigl(\max_i \tilde\delta_i(t)-\min_i \tilde\delta_i(t)\bigr)$ of the
\emph{true} trajectory, an offline diagnostic because the grouping uses ground-truth labels---shows,
at $\alpha{=}0.1$, that the half-width grows toward the stability boundary, from
\mondA{}$^\circ$ for scenarios peaking below $45^\circ$ (group $48/51$ calibration/evaluation,
coverage $0.863$), through $94.2^\circ$ for $45$--$90^\circ$ ($44/52$, $0.923$), to
\mondC{}$^\circ$ above $90^\circ$ ($86/75$, $0.800$): coverage is approximately nominal away from
the boundary and below nominal near it. The band is wide precisely where the stability decision
is most critical; it is a conservative screening envelope, not a tight predictive interval.

The directly calibrated pairwise-difference score $R_{\mathrm{pw}}$ bounds the prediction error of the pairwise angle
difference---the exact quantity in the stability criterion---to \pairq{}$^\circ$ at
$\alpha{=}0.1$ (coverage \paircov{}), below the $\pi$ threshold. Making a decision from this
bound still requires the \emph{predicted} spread to satisfy
$\mathrm{spread} + \hat q_{\mathrm{pw}} < \pi$; on the evaluation half this rule resolves
\decfrac{} of the predicted-stable scenarios at $\alpha{=}0.1$ (rising to $31\%$ at
$\alpha{=}0.3$, $57\%$ at $\alpha{=}0.4$, and $67\%$ at $\alpha{=}0.5$ as $\hat q_{\mathrm{pw}}$
tightens), and the empirical wrong-release rate of the predicted-stable gate is \wrrate{}
($\wrnum{}$ of \wrdem{}), so any deployment must retain a TDS fallback for the unresolved
fraction. [Rev. R4] For deployment, rolling calibration windows refreshed on recent contingencies, and
grouping rules that depend only on observable inputs rather than ground-truth labels, are the
practical counterparts of the offline Mondrian diagnostic; these are left as future work.
Per-machine calibration with Bonferroni/\v{S}id\'ak corrections does not tighten the
bound, and calibrating on the validation set instead of held-out scenarios
yields only $0.61$ coverage at the nominal $0.90$ target: the held-out-bus distribution is
genuinely harder, so calibrating on held-out scenarios is necessary rather than a leakage
concern. The calibration does not transfer across operating points (the $N{-}0$-calibrated interval covers only $0.51$ of
stable $0.8$-load scenarios and $0.72$ at $0.9$ load), consistent with the operating-point
limitation of Section~\ref{sec:discussion}. [Rev. R6] The derived CCT band is vacuous in the reported configuration: with
$2\hat q = 296.2^\circ$ at $\alpha{=}0.1$, the conservative endpoints land at the search lower
bound ($10$ ms) on all $12$ held-out buses, so no scenario is resolved; conformalized monotone
regression over $t_c$ is a candidate for stabilizing the band and is left as future work.
In the gated regime (evaluated on the disjoint half): gate purity \gatepurity{}, coverage
over predicted-stable scenarios \gatecovAll{}, and oracle-conditioned coverage \gatecovTS{}
(\gatenum{} scenarios; the last conditions on ground-truth labels and is not a deployment
guarantee, since conditioning on the prediction breaks exchangeability). [Rev. R1] Alternative scores trade coverage semantics for tighter bands: an
amplitude-normalized (relative) score bounds the error to $2.3\%$ of each scenario's peak
amplitude (coverage $0.854$), a first-swing-window score bounds the first swing to
$70.8^\circ$ (coverage $0.865$; an offline diagnostic, since the window is defined by the true
trajectory), and a time-averaged score yields $47.7^\circ$ (coverage $0.893$) but bounds the
time average rather than the uniform-in-time error; none of these bounds the pairwise spread
directly, so they complement rather than replace $R_{\mathrm{pw}}$. Perturbation-scaled and
input-adaptive scores are candidates for tightening near-boundary bands and are left as future
work, together with adaptive conformal
methods~\cite{gibbs2021adaptive, lei2014distribution} and conformalized quantile
regression~\cite{romano2019conformalized}.

\begin{figure}[t]
\centering
\includegraphics[width=\columnwidth]{fig_conformal.pdf}
\caption{Conformal band validity and adaptivity. Left: empirical coverage (with binomial
intervals) against the target coverage level $1-\alpha$. Right: offline Mondrian diagnostic---the grouping
uses the ground-truth post-fault peak of each scenario---shows the band half-width growing toward
the stability boundary; bars are annotated with group size (calibration/evaluation) and per-group
coverage.}
\label{fig:conformal}
\end{figure}

\subsection{The negative result: input encoding versus residual loss}
The auxiliary physics-residual loss does not help in this setting: with the correctly scaled
residual of \eqref{eq:physloss}, the pooled stable-trajectory RMSE is $39.2^\circ$ at
$\lambda{=}0.001$ (versus \rmseAll{}$^\circ$ for the data-only operator) and $76.5^\circ$ at
$\lambda{=}0.01$ (Table~\ref{tab:ablation})---no improvement at the mildest weight and harmful at
larger weights---and denser collocation does not help either ($48.9^\circ$ pooled, $92.9\%$
accuracy). Without the correct post-fault initial condition the swing-equation residual is a weak
constraint satisfied by many trajectories, so it adds regularization without information. This conclusion is scoped to the present system, protocol, weights, and architecture.

\subsection{Near-boundary sensitivity and architecture ablations}
Training on boundary-aware data (clearing times sampled within $[0.4,1.2]\times$CCT per fault
bus) raises the pooled post-fault RMSE from \rmsePost{}$^\circ$ (uniform) to $70.8^\circ$
(boundary-aware): near-boundary trajectories respond sharply to small parameter changes at the
stability bifurcation, the regime the conformal band is designed to flag. Stratifying the test scenarios by their distance to the TDS CCT quantifies this:
more than $100$ ms inside the boundary the misclassification rate is $8.8\%$; within $30$ ms it
rises to $20.0\%$ ($n{=}60$); and within $10$ ms it reaches $46.2\%$ ($n{=}13$, descriptive). The
sharpest decisions are the least reliable. Replacing the branch network with self-attention does
not improve accuracy ($116.1^\circ$ RMSE, $92.8\%$ accuracy at $p{=}512$, versus
\rmseAll{}$^\circ$ and \pctACC{}), consistent with the severity encoding already capturing the
machine-interaction structure, and a Fourier-neural-operator trunk likewise gives no improvement
($47.4^\circ$ RMSE, $92.6\%$ accuracy). This supports the choice of a
simple DeepONet with a physics-informed encoding.

\subsection{Robustness to measurement noise}
The encoding $a_i(k)=(P_{mi}-P_{ei}^{\mathrm{f}})/M_i$ is itself computed from estimates of
mechanical power, inertia, and fault-on electrical power, which carry measurement error.
Perturbing the encoded acceleration vector with zero-mean Gaussian noise, the pooled stable RMSE
rises from \rmseAll{}$^\circ$ to $115.4^\circ$ and accuracy drops from \pctACC{} to $89.1\%$ at
$5\%$ noise ($152.3^\circ$/$82.9\%$ at $10\%$, $213.6^\circ$/$73.0\%$ at $20\%$). The operator is therefore only moderately robust to
\emph{feature} noise; correlated or biased parameter errors and conformal coverage under noise
remain to be tested.

\subsection{Scalability to the IEEE 118-bus system}
To probe scalability beyond a single test system, we evaluate the operator on a custom dynamic
test case built on the IEEE 118-bus static network: the $54$ machine parameters ($H$, $x'_{d}$)
are taken from the grouped tables of Demetriou et al.~\cite{demetriou2017dynamic} and assigned to
pandapower's case118 generators by rated capacity. No per-machine bus mapping and no per-unit
base conversion are applied ($H$ and $x'_{d}$ are interpreted directly on the system MVA base,
a self-consistent convention); this is a synthetic dynamic case on the IEEE 118-bus network
rather than a reproduction of the Demetriou system; results probe scalability of the
\emph{method}, not of any benchmark. With $24$ fault buses held out entirely from training
($3760$ training and $960$ test scenarios), the retrained operator attains \pctACCb{}
classification accuracy and $121.9^\circ$ stable-trajectory RMSE. The test set is $25.4\%$
stable, so a majority-class baseline would score $74.6\%$; \pctACCb{} therefore reflects real
discrimination, close to the 39-bus \pctACC{} despite $5.4\times$ more generators and $3.0\times$
more buses. The higher trajectory RMSE reflects the larger state space and the much lower CCTs of
the assigned 118-bus case ($0.008$--$0.135$ s), with a CCT error of $9.9$ ms mean absolute (max
$22.3$ ms) among the 11 buses with a defined first crossing. The split-conformal interval attains $0.918$ coverage at
$\alpha{=}0.1$ with a half-width of $596.4^\circ$---larger than the $180^\circ$ stability
threshold, so on the 118-bus case the interval validates the coverage mechanism but is vacuous
for decisions. [Rev. R5] The directly calibrated pairwise
score does not recover a nontrivial resolvable fraction there either: at $\alpha{=}0.1$ its
half-width is $1057.6^\circ$ (coverage $0.918$), and even at $\alpha{=}0.5$ ($91.7^\circ$,
coverage $0.582$) the rule $\mathrm{spread}+\hat q_{\mathrm{pw}}<\pi$ resolves $0$ of the $88$
predicted-stable evaluation scenarios, because the predicted spreads themselves sit near the
$\pi$ threshold; richer calibration data and system-specific scores remain required.

\section{Discussion and Limitations}
\label{sec:discussion}
First, the post-fault trajectory accuracy (\rmsePost{}$^\circ$ pooled RMSE) is limited and
concentrated near the stability boundary, where the classical model response to small parameter
changes is sharpest. Second, the conformal band is calibrated on stable scenarios and covers the
finite output grid: it is a screening band for the stable subpopulation, not a stability
certificate; near-boundary and unstable scenarios are conveyed through the classification, the
wrong-release statistics of Section~\ref{sec:results}, and the heuristic CCT band rather than a
numeric trajectory bound, and on the 118-bus case the per-machine band is currently vacuous. The $\pi$-in-$4$ s criterion is the operational label of this study; whether the first
swing governs multi-swing instability in every scenario of richer models was not established
here. Third, governors, power-system stabilizers, and dynamic loads are not modeled; the one-axis,
exciter, and damped extensions retrain the same architecture rather than transferring weights. Fourth, the pointwise
MLP baseline is slightly worse than the operator on trajectory RMSE at full data
($43.1^\circ$ versus \rmseAll{}$^\circ$ pooled), and the low-data advantage is indicative rather
than established (Section~\ref{sec:sampleeff}); under the \emph{same} conformal and CCT protocol
the MLP is not worse, so the operator's remaining advantages are the lower trajectory RMSE,
continuous-time evaluation, and the low-rank structure, and certification tightness is not one of
them. Fifth, the operator is trained at a fixed operating point and does not generalize to
substantially different load levels: classification accuracy drops to $75\%$ at $20\%$ lower
load, and training across multiple load levels does not resolve this (a model trained on four
load levels attains only $65\%$ on a held-out level). Deployment is therefore a screening tool
\emph{within} its calibrated operating domain; detecting and rejecting out-of-domain operating
points ([Rev. R10] e.g., thresholds on the encoding norm, statistics of the predicted spread,
or conformal-residual monitors) is left as required future work, and
operating-point generalization remains the binding obstacle to broader deployment. Sixth, only bolted three-phase
bus faults are modeled; the encoding is continuous in the fault admittance, so finite-impedance
faults map to nearby encodings, and a fault along a line corridor can in principle be represented
by the admittance of the faulted, split line; neither case has been tested. [Rev. R3] Further
follow-ups include physically informed graph baselines (GNNs fed with $Y$-bus or PTDF/LDF
features) and hybrid schemes that use TEF/IMEAC to refine the operator's CCT near the boundary.

\section{Conclusion}
We presented a physics-encoded neural operator for transient stability assessment that predicts
the full post-fault rotor-angle trajectory and the critical clearing time, with a split-conformal
layer providing a finite-sample, distribution-free trajectory band for the stable subpopulation
on the output grid. The key enabler is a physics-informed input encoding---the swing-equation
acceleration vector, extended by a post-fault severity vector---which lets the operator generalize
to unseen fault locations and to held-out $N{-}1$/$N{-}2$ line-trip topologies. On the IEEE 39-bus
system the operator runs $\approx\timSU{}\times$ faster than time-domain simulation on identical
hardware, classifies stability with \pctACC{} accuracy under the max-over-time $\pi$-separation
criterion, and attains a mean absolute CCT error of \ccterr{} ms against an event-aligned reference; the
directly calibrated pairwise-error bound (\pairq{}$^\circ$ at $\alpha{=}0.1$) is coverage-valid, with a
wrong-release rate and resolvable fraction reported in Section~\ref{sec:results} and a TDS
fallback for the remainder. A negative result shows that, for this system and protocol, the
physics-informed input encoding is more effective than a physics-residual loss, and the encoding
carries over unchanged to one-axis, exciter, and damped generator models and to a synthetic
118-bus dynamic case (\pctACCb{} accuracy), where trajectory-level bands remain vacuous for
decisions. Together with the operating-point limitation, these results define a fast,
uncertainty-aware screening tool for TSA within its calibrated domain, and a concrete
agenda---out-of-domain detection, multi-swing and multi-operating-point calibration, and tighter
near-boundary scores---for its deployment.

\section*{Acknowledgments}
During the preparation of this work the authors used DeepSeek, a generative AI language
model, to polish the language and improve the readability of the manuscript. After using this tool, the
authors reviewed and edited the content as needed and take full responsibility for the content
of the publication.

\section*{Declarations}
\noindent\textbf{Funding.} The authors declare no specific funding for this work.\\
\textbf{Conflict of interest.} The authors declare no competing interests.\\
\textbf{Data availability.} The code reproducing all experiments---scenario generation,
model training for all variants, conformal calibration, CCT binary search, and the sources for
Figures~1--3 and Tables~I--III---is available at
\url{https://github.com/La0Fe1/operator-transient-stability}, together with the random seeds,
software versions, and environment file.\\
\textbf{CRediT authorship contribution statement.} Zhenyu Liu: Conceptualization, Methodology,
Software, Validation, Formal analysis, Investigation, Writing -- original draft, Writing --
review \& editing. Jiale Wang: Validation, Writing -- review \& editing. Jiayong Liu: Validation,
Writing -- review \& editing. Zhicheng Zeng: Validation, Writing -- review \& editing.

\bibliographystyle{elsarticle-num}
\bibliography{references}

\section*{Author biographies}
\noindent\textbf{Zhenyu Liu} is pursuing the B.Eng. degree at the International Energy College,
Jinan University, Guangzhou, China. His research interests include power system stability and
machine learning.\\
\textbf{Jiale Wang} is pursuing the B.Eng. degree at the International Energy College, Jinan
University, Guangzhou, China. His research interests include power system simulation and AI.\\
\textbf{Jiayong Liu} is pursuing the B.Eng. degree at the International Energy College, Jinan
University, Guangzhou, China. His research interests include power system dynamics.\\
\textbf{Zhicheng Zeng} is pursuing the B.Eng. degree at the International Energy College, Jinan
University, Guangzhou, China. His research interests include power system analysis.

\end{document}
